% =============================================================================
% PAPER B — the VERITAS-Bench benchmark/tool paper.
% Target: IEEE INDICON 2026 (deadline 31 Aug 2026) or a Springer CCIS/LNNS venue.
% IEEEtran conference format, max 6 pages.
%
% RELATION TO PAPER A (ethics — read this):
%   Paper A ("A Validity Audit ... for Dysarthric-Speech Severity") reports the
%   EMPIRICAL findings (five confounds, EMA grounding, cycle-consistency,
%   inversion kinematics). THIS paper reports the TOOL: its six controls, the
%   threshold policy, the plug-in API, the report-card schema, and report cards
%   for four metrics — including an audit of the strongest published metric
%   family, which does not appear in Paper A.
%   These are distinct contributions and may be submitted in parallel, BUT you
%   must post Paper A to arXiv FIRST and cite it here (\cite{companion}) so the
%   relationship is transparent. Fill in the arXiv ID before submitting.
%
% BEFORE SUBMITTING: fill \author{}, keep <= 6 pages, add the arXiv ID.
% =============================================================================
\documentclass[conference]{IEEEtran}
\IEEEoverridecommandlockouts

\usepackage{cite}
\usepackage{amsmath,amssymb}
\usepackage{graphicx}
\usepackage{booktabs}
\usepackage{url}

\newcommand{\ns}{\emph{n.s.}}

\begin{document}

\title{VERITAS-Bench: A Validity Test-Suite for\\
Frozen Self-Supervised Clinical-Speech Metrics}

% ---------------------------------------------------------------------------
% TODO: fill in your name, affiliation and email.
% ---------------------------------------------------------------------------
\author{\IEEEauthorblockN{FIRSTNAME LASTNAME}
\IEEEauthorblockA{\textit{Department / Institute} \\
\textit{University Name}\\
City, India \\
your.email@example.com}
}

\maketitle

\begin{abstract}
Training-free clinical-speech metrics built on frozen self-supervised (SSL)
models are proliferating: a severity score is obtained by comparing a patient's
frozen embeddings to healthy-control references, with no labels and no
fine-tuning. Such a metric is typically validated by a single number---a
correlation with clinical severity on one corpus. We argue that this is not
validation, and we provide the missing procedure. \textbf{VERITAS-Bench} takes
\emph{any} per-speaker metric and a corpus and runs six confound controls---
predictive validity, speaker-identity leakage under a leave-one-control-out
reference, silence/pausing, speech energy, cross-corpus transfer, and
representation-level speaker identifiability---emitting a machine-readable
PASS/WARN/FAIL report card against documented thresholds. We report cards for
four metrics on real dysarthric speech (TORGO, $n{=}15$ graded; IPVS $n{=}61$;
EasyCall $n{=}55$), including an independent reimplementation of the strongest
published training-free metric family. Every metric ``passes'' its headline
correlation ($\rho{=}0.74$--$0.77$) yet each fails at least one control. The
battery \emph{discriminates}: it certifies the phonological-subspace design as
silence-, energy- and leakage-robust while exposing a trajectory-distance metric
whose signal is $76\%$ silence; and it surfaces a confound no published metric
reports checking---the frozen representation is almost perfectly
speaker-identifiable (AUC $0.996$--$0.999$), a leakage risk inherited by any
reference-based metric built on it. VERITAS-Bench is released.
\end{abstract}

\begin{IEEEkeywords}
benchmark, validity, dysarthria, self-supervised speech models, clinical speech,
confounds, reproducibility
\end{IEEEkeywords}

% ===========================================================================
\section{Introduction}
% ===========================================================================
Frozen SSL speech representations (WavLM~\cite{chen2022wavlm},
HuBERT~\cite{hsu2021hubert}) encode phonetic and articulatory structure, and a
growing literature exploits this to build \emph{training-free} clinical-speech
metrics---severity scores requiring no dysarthric training data. The evidence
offered for such a metric is almost always the same: a Spearman correlation with
clinical severity, on one corpus.

That number is necessary but nowhere near sufficient. A metric can correlate with
severity because it measures the disorder---or because it measures how long the
speaker pauses, how loudly they speak, or simply \emph{who they are}. These are
not hypothetical: all three occur in metrics currently in use. What the field
lacks is not another metric but a \textbf{standard procedure for deciding whether
a metric is measuring what it claims}.

We contribute that procedure. \textbf{VERITAS-Bench} is a released test-suite
that audits any per-speaker clinical-speech metric against six confound controls
and returns a PASS/WARN/FAIL report card. Our contributions:

\noindent\textbf{(i)} Six validity controls with documented thresholds, covering
predictive validity, speaker-identity leakage, silence, energy, cross-corpus
transfer, and representation-level identifiability (Sec.~\ref{sec:controls}).

\noindent\textbf{(ii)} A metric-agnostic API and a machine-readable report-card
schema, so an author can audit a new metric in a few lines (Sec.~\ref{sec:design}).

\noindent\textbf{(iii)} Report cards for four metrics on real dysarthric speech,
\textbf{including an independent reimplementation of the strongest published
training-free metric family}~\cite{frozenssl2026a,frozenssl2026b}
(Sec.~\ref{sec:cards}). The battery is shown to discriminate---certifying good
design where it is present, not merely failing everything.

\noindent\textbf{(iv)} A systemic finding: the frozen representation on which all
these metrics are built is near-perfectly speaker-identifiable (AUC
$0.996$--$0.999$), which is a leakage risk for \emph{any} reference-based
clinical metric and is not reported as checked by any of them.

The empirical confound findings that motivated these controls are reported in a
companion audit~\cite{companion}; this paper contributes the tool, the threshold
policy, and the report cards.

% ===========================================================================
\section{Related Work}
% ===========================================================================
\textbf{Training-free clinical metrics.} The strongest current example estimates
phonological-feature directions from healthy-control speech in a frozen SSL model
and scores a speaker by the $d'$ separation of their phone embeddings along those
directions, reporting severity correlations across thousands of speakers, many
languages and several aetiologies~\cite{frozenssl2026a,frozenssl2026b}. That work
carefully controls token count and alignment quality. It does \emph{not} report
controlling for speaker-identity leakage or speech energy---two of our six
controls. We audit it in Sec.~\ref{sec:cards}, and it largely holds up.

\textbf{Speaker information in SSL layers.} Early SSL layers encode speaker
identity strongly~\cite{zhu2024wavrx}; we operationalise this as a reference-
construction confound (a control speaker scored against a reference containing its
own data) and as a representation-level AUC probe.

\textbf{What is missing.} Benchmarks in clinical speech rank \emph{systems by
accuracy}. None, to our knowledge, audit \emph{whether a metric's accuracy is
confounded}. VERITAS-Bench is a validity benchmark, not a leaderboard: a metric
does not win by scoring highest, it wins by surviving the controls.

% ===========================================================================
\section{The Six Validity Controls}
\label{sec:controls}
% ===========================================================================
Every threshold is a stated policy choice, documented in one place so a card is
auditable. $\alpha{=}0.05$ throughout; the overall verdict is the \emph{worst
applicable} control.

\begin{enumerate}\itemsep2pt
\item \textbf{Predictive validity.} Within-corpus leave-one-speaker-out Spearman
$\rho$ against clinical severity. PASS $\rho\ge0.50$ and $p<\alpha$.
\item \textbf{Speaker-identity leakage (LOCO).} For reference-based metrics,
$\Delta = \rho(\text{naive reference}) - \rho(\text{leave-one-control-out
reference})$. A reference (or a fitted direction) that contains the test speaker's
own data inflates the score. PASS $\Delta\le0.05$.
\item \textbf{Silence confound.} Retention $=\rho(\text{speech-only}) /
\rho(\text{full})$ under MFA~\cite{mcauliffe2017mfa} alignment, plus the silence
frame fraction. PASS if retention $\ge0.70$ \emph{and} silence $<50\%$.
\item \textbf{Energy confound.} Partial Spearman controlling for mean speech
log-RMS energy, relative to the raw $\rho$. PASS if $\ge70\%$ retained.
\item \textbf{Cross-corpus transfer.} Train on the other corpora, test held out.
PASS $\rho\ge0.50$, $p<\alpha$.
\item \textbf{Representation identifiability.} Macro one-vs-rest speaker-
classification AUC of the representation the metric uses. PASS AUC $\le0.70$
(chance $0.5$).
\end{enumerate}

% ===========================================================================
\section{Design and API}
\label{sec:design}
% ===========================================================================
A metric implements one required method, \texttt{predictive\_scores(views)},
returning $\{\text{speaker}\rightarrow\text{score}\}$ with \emph{higher $=$ more
severe}; the battery verifies speaker-disjointness and fails loudly otherwise.
Two optional overrides unlock further controls: \texttt{loco\_scores} enables the
leakage test (reference-based metrics), and a \texttt{has\_representation} flag
exposes the metric's features to the transfer and identity probes. Corpus loaders
yield per-speaker \emph{views} carrying pooled features, silence masks and energy,
so the same metric can be audited on any corpus without change.

The output is a plain JSON report card (one per metric $\times$ corpus): per
control an outcome, a one-line headline with the load-bearing number, and the raw
statistics; plus the threshold policy that produced the verdicts. Cards are
diff-able in version control, so a metric's validity can be tracked as it evolves.

% ===========================================================================
\section{Report Cards}
\label{sec:cards}
% ===========================================================================
We audit four metrics on TORGO~\cite{rudzicz2012torgo} (15 speakers, graded
severity), with IPVS~\cite{dimauro2017ipvs} (61) and
EasyCall~\cite{turrisi2021easycall} (55) as cross-corpus references. All use frozen
WavLM-Large; seed 42; speaker-disjoint throughout. The metrics are: a
trajectory-distance reference metric at a mid layer (\textbf{ATD-L9}) and an early
layer (\textbf{ATD-L1}); a fairly-tuned pooled-embedding Ridge \textbf{probe}; and
\textbf{phono-subspace}, our independent reimplementation, \emph{from the published
description and not the authors' code}, of the phonological-subspace $d'$
metric~\cite{frozenssl2026a,frozenssl2026b} (phone-level embeddings via MFA;
contrast directions for nasality, voicing, stridency, sonorance and continuancy
estimated on healthy controls only; score $=-\overline{d'}$).

\begin{table}[t]
\centering
\caption{VERITAS-Bench report cards on TORGO ($n{=}15$, real, seed 42). Every
metric passes predictive validity, yet each fails at least one control.}
\label{tab:cards}
\small
\begin{tabular}{@{}lcccc@{}}
\toprule
Control & ATD-L9 & ATD-L1 & Probe & \textbf{Phono-sub.}\\
\midrule
predictive validity     & PASS $.77$ & PASS $.69$ & PASS $.74$ & PASS $.74$\\
leakage (LOCO $\Delta$) & PASS $.00$ & \textbf{WARN} $.09$ & N/A & PASS $.00$\\
silence (kept)          & \textbf{FAIL} $24\%$ & --- & WARN $100\%$ & WARN $100\%$\\
energy (kept)           & PASS $96\%$ & --- & PASS $80\%$ & PASS $88\%$\\
cross-corpus            & \textbf{FAIL} \ns & --- & \textbf{FAIL} \ns & \textbf{FAIL} \ns\\
identity AUC            & \textbf{FAIL} $.996$ & --- & \textbf{FAIL} $.996$ & \textbf{FAIL} $.999$\\
\midrule
\textbf{overall}        & \textbf{FAIL} & \textbf{WARN} & \textbf{FAIL} & \textbf{FAIL}\\
\bottomrule
\end{tabular}
\end{table}

\textbf{The battery discriminates.} This is the central result. All four metrics
report a respectable headline correlation ($\rho{=}0.74$--$0.77$) and would, on the
usual evidence, be declared successful. The card separates them:

\emph{Silence.} The trajectory-distance metric loses $76\%$ of its severity signal
when silence is removed (retention $24\%$, FAIL): it is substantially measuring
\emph{pausing}. The phonological-subspace metric retains $100\%$---it is
silence-free by construction, since it only reads frames inside aligned phone
intervals. A benchmark that returned FAIL for everything would be useless; this one
certifies a good design where the design is good.

\emph{Leakage.} The early-layer reference metric is inflated by $\Delta{=}0.089$
under a naive reference (WARN). The phonological-subspace metric's contrast
directions are fitted on the control pool, so a control speaker is scored against
directions containing its own data; re-fitting the directions without that speaker
shifts its raw score, but near-uniformly across controls, so the severity
\emph{ranking} survives ($\Delta{=}0.000$, PASS). This is a genuine robustness
result for that design, obtained by testing a confound its authors do not report
checking.

\emph{Energy.} All metrics retain $80$--$96\%$ of their correlation after
partialling out mean speech energy (PASS).

\emph{Identity.} Every representation is near-perfectly speaker-identifiable (AUC
$0.996$--$0.999$). Because this holds \emph{across all four metrics}, it is not a
defect of any one of them but a \textbf{systemic property of the frozen SSL
representation}: any reference-based clinical metric built on these features
inherits a strong identity-leakage risk, whether or not it manifests as rank
inflation on a given corpus. We are not aware of a published training-free clinical
metric that reports measuring this.

\emph{Cross-corpus.} All four fail to transfer significantly to the graded corpus
($\rho{=}0.43$--$0.47$, $p{\approx}0.08$--$0.11$). We state plainly that this is an
\textbf{underpowered null}: at $n{=}15$ the critical value is $\rho{\approx}0.52$.
The phonological-subspace work demonstrates transfer at
$n{\approx}3{,}374$~\cite{frozenssl2026b}; our battery reports what the available
graded data can support, and labels the limit rather than overclaiming.

% ===========================================================================
\section{Discussion and Limitations}
% ===========================================================================
The report card is deliberately unflattering: no metric we tested, including the
best published one, passes every control. That is the intended function. A single
headline $\rho$ cannot distinguish a metric that measures articulation from one that
measures pausing, loudness, or identity, and the community currently has no standard
way to tell them apart. Six cheap controls do.

\textbf{Limitations.} Graded severity is available for only $n{=}15$ speakers, so the
cross-corpus control is underpowered and its FAIL verdicts must be read as ``not
demonstrated,'' not ``disproved.'' Our phonological-subspace metric is a
reimplementation and may differ from the authors' in detail; we therefore report it
as evidence about the \emph{design}, not as a replication of their numbers. Thresholds
are policy choices, stated explicitly and changeable in one place.

% ===========================================================================
\section{Conclusion}
% ===========================================================================
Frozen-SSL clinical-speech metrics are easy to build, easy to validate badly, and
easy to fool. We release VERITAS-Bench so that a claimed clinical metric can be
audited---by its authors, its reviewers, or its users---against the confounds that a
single correlation cannot exclude. Code, thresholds, and all report cards are
released for reproduction.

\bibliographystyle{IEEEtran}
\bibliography{refs}

\end{document}
